\capitulo{5}{Aspectos relevantes del desarrollo del proyecto}

Desde la propuesta inicial hasta la publicación de \textit{Manualito}, el
trabajo combinó la investigación técnica con el desarrollo de la aplicación.
La organización del proyecto, las decisiones de implementación y las pruebas
realizadas permiten entender cómo se llegó al resultado final.

\section{Inicio del proyecto}
\label{aspectos:inicio}

\textit{Manualito} nace de una propuesta de HP SCDS dentro de su
Observatorio Tecnológico, una iniciativa que conecta a la empresa con las
universidades mediante propuestas de trabajos de fin de
grado~\cite{hp_observatorio_tecnologico}.

Tras una primera valoración de la propuesta, en la reunión inicial se
explicó la idea con más detalle y se debatió la orientación del proyecto.
Las primeras semanas se dedicaron a investigar las tecnologías necesarias
y a realizar pruebas técnicas para comprobar la viabilidad del proyecto.

\section{Metodología de trabajo}
\label{aspectos:metodologia}

El trabajo se organizó mediante una metodología ágil adaptada al desarrollo
individual, con iteraciones de aproximadamente dos semanas. En las reuniones
de seguimiento se revisaban los avances y las dificultades encontradas para
ajustar las prioridades y acordar los objetivos y nuevos requisitos de la
siguiente iteración.

Las tareas se registraron en un tablero Kanban, donde se distinguía el
trabajo pendiente, en curso y terminado. Las tareas pueden consultarse
en el \href{https://github.com/ibaimoya/manualito/milestones?state=closed}{repositorio del
proyecto}, mientras que el Anexo A recoge las estimaciones, el tiempo
dedicado y las desviaciones de cada iteración.

Durante el desarrollo se aplicó un enfoque de desarrollo guiado por
pruebas (\acs{TDD}, del inglés \textit{Test-Driven Development}).\acused{TDD}
La inteligencia artificial ayudó a elaborar las pruebas y el código a partir
de instrucciones concretas, y sus propuestas se revisaron antes de
incorporarlas al proyecto, como se detalla en el apartado~\ref{aspectos:uso-ia}.

\section{Organización del desarrollo}
\label{aspectos:organizacion}

Las \textit{issues} se clasifican mediante etiquetas de tipo
(\texttt{type::...}), prioridad (\texttt{priority::...}) y tema
(\texttt{topic:...}), para distinguir qué se solicita, su urgencia y las
áreas afectadas. Cada una debe tener un solo tipo, una sola prioridad y
uno o varios temas. Las plantillas asignan el tipo y el estado
\texttt{status:needs-triage}, que indica que la clasificación está pendiente
de revisión. Este estado se retira una vez completadas las etiquetas,
mientras que los demás estados \texttt{status:...} son temporales y solo
puede haber uno activo a la vez~\cite{manualito_contributing}.

Con motivo de mantener cierta trazabilidad, en \textit{Manualito},
cada tarea debe estar vinculada a una \textit{issue} y se desarrolla en una
rama creada a partir de \texttt{development}, de modo que sus cambios se
preparan por separado antes de incorporarlos a la aplicación. Para añadir
funcionalidades se utiliza el formato \texttt{feat/<número>-<descripción>},
donde el número corresponde al identificador de la \textit{issue} y la descripción se escribe
en minúsculas, con guiones entre palabras.

En la figura~\ref{fig:aspectos-ramas} se representa el ciclo de vida de un cambio en
este proyecto, tal y como se muestra, la rama de tarea se fusiona con la rama
\texttt{development} cuando está terminada y probada. Desde allí, el conjunto de cambios
que se quiere probar pasa a la rama\texttt{staging}, desde donde se prueba la
aplicación en otro equipo diferente al de desarrollo.
Si la revisión es satisfactoria, se integra en la rama \texttt{master}, que conserva
la última versión estable. Esta estrategia de ramas por entorno, o
\textit{branch per environment}, permite seguir incorporando mejoras en
\texttt{development} mientras se comprueba en \texttt{staging} la versión que
se prepara para publicar~\cite{gitlab_branch_per_environment}. Las \textit{pull
requests} deben enumerar los cambios y mostar los resultados de las pruebas en su
descripción. Las comprobaciones automáticas y los procesos de publicación
asociados a cada evento y rama se detallan en el
apartado~\ref{aspectos:automatizacion}.

\begin{figure}[htbp]
    \centering
    \includegraphics[width=\textwidth]{memoria/aspectos-relevantes-del-desarrollo-del-proyecto/flujo-ramas.pdf}
    \caption{Recorrido del ciclo de vida de una tarea entre las ramas de \textit{Manualito}}
    \label{fig:aspectos-ramas}
\end{figure}
\FloatBarrier

Antes de crear un \textit{commit}, la configuración de \textit{pre-commit}
ejecuta comprobaciones sobre los archivos implicados. Las comprobaciones configuradas  son las siguientes~\cite{precommit_documentacion}:

\begin{itemize}
    \item \textbf{Formato de los archivos.} Elimina los espacios sobrantes al
    final de las líneas, conservando los dos espacios que indican un salto
    de línea en Markdown, y
    ajusta los archivos no vacíos para que terminen con un único salto de
    línea. Estas correcciones excluyen \path{uv.lock} y
    \path{frontend/pnpm-lock.yaml}~\cite{precommit_hooks}.

    \item \textbf{Conflictos pendientes.} Detecta marcas de conflictos sin
    resolver, incluso cuando no hay una fusión en curso.

    \item \textbf{Rama de trabajo.} Impide los \textit{commits} directos en
    \texttt{development}, \texttt{staging} y \texttt{master}, y los realizados
    sin una rama activa. Comprueba que el nombre de la rama combine el tipo
    de cambio, el número de la tarea y una descripción breve en minúsculas,
    con guiones entre palabras. Por ejemplo,
    \texttt{feat/42-mejorar-buscador} representaría una nueva funcionalidad
    (\texttt{feat}) vinculada a la tarea 42 para mejorar el buscador.

    \item \textbf{Formato del mensaje.} Comprueba el cumplimiento del formato de
    \textit{Conventional Commits}~\cite{conventional_commits}.
    Las reglas del proyecto limitan los tipos admitidos y exigen minúsculas
    en el tipo y en el ámbito opcional, con guiones si este contiene varias palabras.
    La descripción debe comenzar en minúscula, tener al menos 10 caracteres
    y no contener emojis, espacios en sus extremos ni un punto final. El
    título completo, incluida la referencia a la tarea, no puede superar
    los 72 caracteres. También reduce los espacios consecutivos del título
    a uno solo.

    \item \textbf{Referencias a tareas.} Comprueba que el título del
    \textit{commit} incluya el número de la tarea de su rama. Lo añade si no
    hay referencias y rechaza el mensaje si solo aparecen otras tareas.
    Las referencias adicionales se agrupan al final del mensaje en una línea
    \texttt{Refs:}, separada del texto por una línea en blanco. Rechaza
    números no válidos, referencias repetidas o menciones como
    \texttt{\#42} fuera de esos lugares, aunque permite enlaces completos
    en el resto del mensaje.

    \item \textbf{Cambios incompatibles.} Cuando el mensaje declara uno,
    exige tanto el signo \texttt{!} en el título como un único pie
    \texttt{BREAKING CHANGE:} acompañado de una explicación.
\end{itemize}

Se excluyen de esta comprobación los mensajes de \texttt{merge} y
\texttt{revert}, así como los que comienzan por \texttt{fixup!},
\texttt{squash!} o \texttt{amend!}.
Si estas correcciones modifican archivos, indica al desarrollador que hay revisarlos y volver a añadirlos al \textit{commit}
antes de completarlo.

La elección entre las plataformas de alojamiento se recoge en el
apartado~\ref{herramientas:organizacion}.
Respecto a la ubicación del repositorio, el proyecto se alojó inicialmente en la instancia privada de GitLab
de HP SCDS, sin embargo, al comienzo del octavo \textit{sprint}, se decidió migrar
el proyecto a GitHub cuando los \textit{runners} de \ac{CI} alojados en dicha
infraestructura empezaron a fallar, lo que dificultaba mantener y mejorar
la integración continua y el despliegue continuo. La adaptación del flujo de
trabajo y la réplica
de las tareas de GitLab se recogen en la
\href{https://github.com/ibaimoya/manualito/issues/63}{\textit{issue} \#63}.
Pese a esto, GitLab conservó una copia del código y su configuración de
integración continua, además del seguimiento del proyecto. Las \textit{issues}
se replicaron en ambas plataformas para mantener su trazabilidad.

\section{Investigación y decisiones técnicas}
\label{aspectos:investigacion}

La elección de motores y modelos no dependía únicamente de la calidad de
sus resultados, ya que también debían ejecutarse con los recursos del equipo
disponible. Para valorar ese equilibrio se prepararon comparaciones del
reconocimiento, la recuperación y la generación de respuestas. Sus resultados
sirvieron de apoyo a las decisiones, junto con las comprobaciones realizadas
durante el uso de la aplicación.

Los ocho estudios principales se encuentran en
\path{docs/benchmarks/}, junto con dos comparaciones anteriores de
preprocesado y modelos de lenguaje~\cite{manualito_benchmarks}.
Los \textit{notebooks} incluyen el método y las mediciones guardadas, mientras
que las decisiones y los resultados más relevantes se explican aquí.

\subsection{Reconocimiento y corrección del texto}

Se compararon Tesseract, PaddleOCR y EasyOCR con imágenes sintéticas y tres
documentos reales. En estos últimos, PaddleOCR obtuvo una tasa media de
error de caracteres del 10,92\,\%, frente al 20,55\,\% de Tesseract.
La variante con \acs{GPU} mantuvo el error de la variante con \acs{CPU},
pero redujo el tiempo medio de 116,6 a 11,4 segundos por imagen.
Tesseract fue más rápido, con 6,9 segundos, aunque a costa de cometer más
errores. La figura~\ref{fig:benchmark-ocr} resume esta comparación, en la que
el error se calcula sobre fragmentos transcritos y cada imagen recibe el
mismo peso~\cite{manualito_benchmarks}.

\begin{figure}[htbp]
    \centering
    \includegraphics[width=\textwidth]{memoria/aspectos-relevantes-del-desarrollo-del-proyecto/benchmarks/ocr.pdf}
    \caption{Error de caracteres y tiempo de los motores de \acs{OCR}}
    \label{fig:benchmark-ocr}
\end{figure}

Se priorizó PaddleOCR con \acs{GPU} cuando había memoria gráfica suficiente,
por su combinación de calidad y tiempo. Sin embargo, reservar parte de esa
memoria al reconocimiento puede dejar menos espacio para un modelo de
lenguaje grande y, si el modelo no cabe por completo, parte de su ejecución
puede trasladarse al procesador y a la memoria del sistema, lo que se conoce
como \textit{offloading} a \acs{CPU} y, en conseciencia, reducir la
velocidad~\cite{ollama_contexto,ollama_memoria}. Por ello se conservaron
alternativas de \acs{OCR} que permiten reservar la tarjeta gráfica al \ac{LLM}.
La muestra real era pequeña y Tesseract obtuvo menos error en las imágenes
sintéticas, por lo que no se consideró que un motor fuera superior en todos
los documentos.

El preprocesado se evaluó por motor y frente a la imagen sin modificar.
Sobre diez imágenes, se eligieron la ampliación de 2,5 veces y la conversión
a grises para Tesseract, y el ajuste de contraste local para PaddleOCR.
No bastaba con reducir el error medio, también se exigía mejorar al menos
la mitad de las imágenes. Estas opciones se incorporaron al servicio, sin
asumir que aplicar más transformaciones mejoraría cualquier documento.

La corrección posterior con ayuda del \ac{LLM}, se planteó como una ayuda conservadora. Su prioridad
era no alterar reglas que ya estuvieran bien y no añadir información falsa, antes que corregir el mayor
número posible de líneas. Por ello, el modelo solo interviene en una franja
de confianza del reconocimiento y sus propuestas pasan por un consenso.
Este consenso consiste en solicitar al mismo modelo tres propuestas de corrección para cada línea
cuya puntuación de confianza esté dentro del rango establecido y solo aceptar aquellos cambios que
coinciden en al menos dos de las propuestas.
En el conjunto real se corrigieron completamente 2 de las 24 líneas con
errores evaluadas, sin modificar las 12 líneas correctas. El error de
caracteres bajó del 18,49\,\% al 18,04\,\% sobre las mismas líneas.

El corrector aislado resolvió 31 de 34 líneas sintéticas que tenían errores, pero
esa cifra no representa el resultado del flujo completo, que trabaja con
otros casos y aplica los filtros anteriores. Se aceptó una mejora reducida
como complemento del reconocimiento, teniendo en cuenta que, la fuente de verdad sigue siendo el manual,
porque una propuesta del modelo puede parecer razonable y, aun así,
introducir información que no estaba en el documento.

\subsection{Fragmentación y recuperación}

La comparación de \textit{embeddings} utilizó 36 pasajes sintéticos y
24 consultas. E5-small recuperó un pasaje válido entre los tres primeros
en 23 consultas, igual que E5-base y MiniLM multilingüe, pero lo situó
primero con mayor frecuencia. Además, necesitó 1,92 milisegundos por
documento frente a los 7,24 de E5-base. Esto apoyó su elección para la
ejecución local, aunque solo se conservan los resultados agregados y el
conjunto no permite asegurar el mismo rendimiento con cualquier manual.

En la fragmentación se compararon siete configuraciones sobre tres manuales
sintéticos y diez consultas, entre las que destacó la división recursiva de
1024 caracteres con un solapamiento de 128 utilizada en la aplicación.
Esta configuración recuperó el fragmento esperado entre los tres primeros
resultados en las diez consultas, con un tiempo medio de 8,51 milisegundos
por consulta.

La continuidad del texto se evaluó mediante una medida basada únicamente
en el signo de puntuación al final de cada fragmento, por lo que el
porcentaje obtenido no permite distinguir de forma fiable entre una frase
interrumpida y un corte válido.

La recuperación se comprobó también en el servicio real, con 24 preguntas
sobre Catán y Monopoly que incluían fallos detectados durante el desarrollo.
La figura~\ref{fig:benchmark-recuperacion} recoge las mediciones
sucesivas~\cite{manualito_benchmarks}, cuyas fechas identifican las ejecuciones
guardadas y no pruebas independientes de cada cambio.
El fragmento esperado pasó de aparecer entre los cinco primeros en
19 preguntas a hacerlo en 22, y en la última medición se encontró entre
los diez primeros en las 24. Esto permite observar la evolución del buscador
en el conjunto evaluado, pero encontrar el texto adecuado no demuestra
todavía que la respuesta generada a partir de él sea correcta.

\begin{figure}[htbp]
    \centering
    \includegraphics[width=\textwidth]{memoria/aspectos-relevantes-del-desarrollo-del-proyecto/benchmarks/recuperacion.pdf}
    \caption{Evolución de la recuperación}
    \label{fig:benchmark-recuperacion}
\end{figure}

\subsection{Modelos de lenguaje y recursos disponibles}

Las comparaciones de modelos se utilizaron como orientación, no como una
selección automática. Se valoraron respuestas a preguntas directas y a
casos con información insuficiente, contradicciones o excepciones, junto
con la velocidad y la memoria gráfica ocupada. Las pruebas se complementaron
con simulaciones de uso real de la aplicación para valorar el comportamiento
de los modelos fuera de los casos preparados.

La comparación evaluó ocho modelos con 24 casos y tres
repeticiones en una NVIDIA GeForce RTX 4070 SUPER de 12 GB.
Gemma 4 E4B obtuvo una puntuación de fidelidad de 0,972 sobre 1 en los
casos adversariales, según los jueces automáticos y, con un consumo de
9,84 GB de memoria gráfica.

Por otr parte en el estudio de diez modelos pequeños,
Granite 3.3 2B obtuvo 0,944 y consumiendo alrededor de 1,96 GB.
Cada punto de la figura~\ref{fig:benchmark-modelos} representa un modelo según
su puntuación media de fidelidad en los casos adversariales y la memoria
gráfica registrada~\cite{manualito_benchmarks}. Los resultados apoyaron la
elección de Gemma y Granite para los perfiles alto y ligero, respectivamente,
aunque proceden de estudios con distintos conjuntos de jueces y no deben
leerse como una clasificación conjunta.

\begin{figure}[htbp]
    \centering
    \includegraphics[width=\textwidth]{memoria/aspectos-relevantes-del-desarrollo-del-proyecto/benchmarks/modelos.pdf}
    \caption{Fidelidad media y memoria gráfica usada de los modelos}
    \label{fig:benchmark-modelos}
\end{figure}

Aparecieron desacuerdos entre los jueces, lo que limitó el peso de estas puntuaciones en la
decisión final. Por ello, se realizaron una serie de simulaciones de uso que aportaron una comprobación complementaria. Del mismo modo, el
límite de 4 GB del estudio ligero se contrastó con la memoria ocupada en
una tarjeta de mayor capacidad, no mediante una prueba física en una tarjeta
de 4 GB. La elección de un perfil depende, por tanto, de los recursos que
queden disponibles para el conjunto de la aplicación.
\FloatBarrier

\section{Seguridad}
\label{aspectos:seguridad}

Para proteger las cuentas, evitar el acceso a datos ajenos y limitar el uso
abusivo del servicio, se incorporaron las siguientes medidas de seguridad:

\begin{itemize}
    \item \textbf{Credenciales.} Las contraseñas se guardan mediante
    Argon2id y se comprueba su longitud~\cite{owasp_password_storage}.
    Si la cuenta no existe, se utiliza un \textit{hash dummy} para comprobar
    igualmente la contraseña y reducir las diferencias
    de tiempo que podrían revelar si existe. El inicio de sesión
    devuelve el mismo error ante una cuenta inexistente o una contraseña
    incorrecta, mientras
    que la recuperación y el reenvío de verificación utilizan mensajes
    que no confirman si el correo está registrado~\cite{owasp_authentication}.

    \item \textbf{Sesiones y recuperación.} Las sesiones utilizan
    identificadores aleatorios cuyos \textit{hashes} se guardan en la base
    de datos. La \textit{cookie} de sesión incorpora \texttt{HttpOnly},
    \texttt{Secure} y \texttt{SameSite=Lax} para impedir su lectura desde
    JavaScript, exigir conexiones cifradas y limitar su envío desde otros
    sitios web. En cada petición que requiere autenticación, El servidor
    comprueba que la sesión no haya caducado y que la cuenta
    siga activa~\cite{owasp_session_management}. Por otro lado, los enlaces de verificación
    y recuperación caducan y son de un solo uso. Además, cambiar la contraseña
    exige indicar la anterior y revoca las otras sesiones, mientras que
    restablecerla las revoca todas.

    \item \textbf{Peticiones y permisos.} Los ataques de \ac{CSRF} intentan
    provocar acciones desde otra web aprovechando una sesión abierta.
    Para protegerse, las operaciones autenticadas que modifican datos
    exigen un \textit{token} que el servidor comprueba contra el asociado
    a la sesión~\cite{owasp_csrf}. También se comprueban la propiedad y
    visibilidad de los recursos y los permisos de administración, para
    impedir el acceso a datos u operaciones no autorizados.

    \item \textbf{Límites y archivos.} Se limita la frecuencia de las
    peticiones a rutas sensibles y el número de cálculos de contraseñas
    simultáneos para dificultar los intentos repetidos de acceso y la
    saturación del servicio. Frente a \ac{ReDoS}, donde ciertos textos
    pueden hacer que una expresión regular tarde excesivamente, hasta en el
    peor caso llegar a poder tirar la aplicación, se
    han sustituido patrones problemáticos de la validación del correo y del
    procesamiento del texto por operaciones cuyo tiempo crece de forma
    proporcional al tamaño de la entrada~\cite{owasp_input_validation}.
    En las subidas se comprueban el formato real,
    el tamaño, el número de páginas y las dimensiones de las imágenes.
    También se rechazan rutas que salgan del directorio de almacenamiento
    y enlaces simbólicos~\cite{owasp_file_upload}.

    \item \textbf{Contenido en el navegador.} Los ataques de \ac{XSS}
    buscan ejecutar código ajeno en la página~\cite{mdn_xss}. Para reducir
    este riesgo, las páginas de los \acs{PDF} se muestran como imágenes,
    que se sirven con su tipo de contenido y la cabecera
    \path{X-Content-Type-Options} con el valor \texttt{nosniff} para que el navegador respete
    ese tipo~\cite{mdn_nosniff}. Las respuestas se muestran con
    \texttt{react-markdown}, que muestra las etiquetas como texto y bloquea los enlaces
    que utilizan protocolos no permitidos, como \texttt{javascript:}~\cite{react_markdown_security}.

    \item \textbf{Inyección de instrucciones (\textit{prompt injection}).}
    Un mensaje o un manual puede contener instrucciones que intenten desviar
    al modelo de su tarea~\cite{anthropic_prompt_injection}. Para reducir
    este riesgo, el \textit{prompt} de respuesta distingue las instrucciones
    internas del historial, los fragmentos del manual y la pregunta.
    Al modelo se le indica que utilice esos bloques como datos, no como
    nuevas instrucciones, y que rechace las peticiones de revelar o modificar
    sus reglas internas. Sin embargo, estas indicaciones no garantizan que todos los
    intentos de manipulación se rechacen. Esto es una limitación conocida,
    debido a que implantar guardarraíles es demasiado laborioso para el
    poco riesgo que tiene este tipo de ataque en este entorno. Es por ello que,
    se trata de minimizar las consecuencias de este tipo de ataque mediante un
    filtrado el previo del contexto que se le da según los permisos del usuario que consulta.

    \item \textbf{Bastionado de imágenes y contenedores.} Las imágenes de
    los servicios propios se configuran para ejecutarse con un usuario sin
    privilegios de administrador. Para estos servicios, Docker Compose
    retira las capacidades de Linux, se borran gestores de paquetes
    y herramientas innecesarias, y se activa \path{no-new-privileges}, que
    impide obtener privilegios adicionales al ejecutar otros
    programas~\cite{docker_container_run_seguridad}. El sistema de archivos
    se monta en modo de solo lectura, salvo los volúmenes y directorios
    temporales necesarios~\cite{docker_compose_servicios}. Los contenedores
    de inicialización conservan los permisos que necesitan para preparar
    esos volúmenes. Con ello se busca limitar el alcance de un posible ataque.

    \item \textbf{Despliegue y registros.} Caddy cifra el acceso y aplica
    \path{Content-Security-Policy} para restringir los recursos que puede
    cargar el navegador~\cite{owasp_csp}. Las bases de datos y los servicios
    de procesamiento no se exponen directamente a Internet, solamente a la
    red interna que existe por detrás de Caddy. Las
    credenciales de infraestructura se suministran como secretos. Los
    \textit{logs} filtran campos sensibles para evitar exponerlos y Caddy oculta las
    \textit{cookies} y credenciales de acceso en sus registros.
\end{itemize}

El Anexo C amplía el diseño de seguridad y la configuración del despliegue.

\section{Pruebas}
\label{aspectos:pruebas}

Los \textit{benchmarks} del apartado~\ref{aspectos:investigacion} comparan
alternativas de procesamiento. Las pruebas del software comprueban, en cambio,
que sus componentes se comporten como se espera, también al introducir cambios.

El proyecto incluye\textbf{ pruebas unitarias} y de \textbf{integración}, ejecutadas con
pytest en el servidor y Vitest en la interfaz. Las unitarias comprueban
operaciones concretas, como el orden de los resultados al combinar
búsquedas, mientras que las de integración verifican que varios componentes
funcionen correctamente juntos. Por ejemplo, al crear un manual se
comprueba que sus datos puedan consultarse en PostgreSQL y que sus imágenes
estén disponibles en el almacenamiento.

Los informes de cobertura permiten identificar qué partes del código se
ejecutan durante las pruebas y orientar la revisión de las que quedan fuera.
Esta medida se interpreta junto con los casos comprobados, porque recorrer
una línea no garantiza haber evaluado todos sus comportamientos. La
generación de estos informes y su análisis con SonarQube se explican en el
apartado~\ref{aspectos:automatizacion}.

Para comprobar las funcionalidades desde la interfaz, también se realizaron
\textbf{pruebas de aceptación} mediante \textbf{agentes de inteligencia artificial} que
utilizaron Playwright para interactuar con el
navegador~\cite{playwright_library}. Los agentes dirigían las acciones sobre
la aplicación y comprobaban su respuesta, con el propósito de contrastar
el comportamiento observado con los requisitos del proyecto.

El Anexo D recoge la organización y las instrucciones de ejecución de las
pruebas automatizadas, mientras que la participación de la inteligencia
artificial se amplía en el apartado~\ref{aspectos:uso-ia}.

\section{Integración continua y despliegue continuo}
\label{aspectos:automatizacion}

La integración continua actúa como malla de seguridad al repetir ciertas
comprobaciones al enviar cambios y al abrir \textit{pull requests} en GitHub
o \textit{merge requests} en GitLab~\cite{github_ci,gitlab_ci}.
\acused{CI} En \textit{Manualito}, estas comprobaciones están replicadas en
GitHub Actions y GitLab, aunque el repositorio de referencia se encuentra
en GitHub como ya se ha mencionado anteriormente. Cada repositorio conserva su configuración, que
agrupa el trabajo en los siguientes bloques~\cite{manualito_workflows,
manualito_gitlab_ci}:

\begin{itemize}
    \item \textbf{Servidor.} Ruff revisa el código, mypy comprueba los tipos
    y pytest ejecuta las pruebas y obtiene la cobertura, con PostgreSQL
    disponible para las pruebas que necesitan una base de datos.
    \item \textbf{Migraciones.} Alembic aplica los cambios sobre una base
    de datos y comprueba que los modelos no requieran una migración pendiente.
    \item \textbf{Interfaz.} Se revisan el código con ESLint y los tipos
    con TypeScript, se ejecutan las pruebas con Vitest y se genera la
    versión de producción.
    \item \textbf{Despliegue.} Se valida la configuración de Caddy,
    Docker Compose y Terraform antes de utilizarla para arrancar los servicios.
\end{itemize}

La figura~\ref{fig:aspectos-cicd} muestra los trabajos que pueden ejecutarse
en paralelo y cuáles deben completarse correctamente antes del análisis
de SonarQube.

\begin{figure}[htbp]
    \centering
    \includegraphics[width=\textwidth]{memoria/aspectos-relevantes-del-desarrollo-del-proyecto/flujo-cicd.pdf}
    \caption{Flujo de integración continua en GitHub Actions}
    \label{fig:aspectos-cicd}
\end{figure}

También se auditan las dependencias de producción de la interfaz y se muestran sus
avisos de vulnerabilidades altas o críticas. SonarQube Cloud reúne los informes de
cobertura y analiza problemas de calidad, duplicación y complejidad. Su
análisis espera el resultado del control de calidad y falla si no se supera.

La tabla~\ref{tabla:aspectos-ci-ramas} resume las comprobaciones por evento
en ambas plataformas. En GitHub, las \textit{pull requests} activan el flujo
cuando se dirigen a \texttt{development}, \texttt{staging} o \texttt{master},
mientras que GitLab lo aplica a las \textit{merge requests} y evita duplicar
la ejecución de rama cuando ya existe una solicitud abierta.

\begin{table}[htbp]
    \centering
    \small
    \begin{tabular}{@{}p{0.53\textwidth}p{0.43\textwidth}@{}}
        \toprule
        \textbf{Evento y rama} & \textbf{Comprobaciones} \\
        \midrule
        \textit{Commit} a una rama de tarea & Controles locales al crear el
        \textit{commit}. La \acs{CI} se activa con su solicitud de integración. \\
        \addlinespace
        \textit{Pull request} o \textit{Merge request} & Comprobaciones comunes y SonarQube Cloud. \\
        \addlinespace
        \textit{Commit} a \texttt{development} o \texttt{staging} & Comprobaciones comunes. \\
        \addlinespace
        \textit{Commit} a \texttt{master} & Comprobaciones comunes y SonarQube Cloud. \\
        \bottomrule
    \end{tabular}
    \caption{Comprobaciones configuradas por evento y rama en GitHub y GitLab.
    Elaboración propia a partir de los flujos de ambos repositorios.}
    \label{tabla:aspectos-ci-ramas}
\end{table}
\FloatBarrier

El flujo también puede iniciarse manualmente sobre una rama elegida, con
SonarQube si se selecciona \texttt{master}. La prueba en otro equipo al
pasar a \texttt{staging} complementa estos controles automáticos. El
Anexo D recoge las instrucciones para ejecutar las comprobaciones.

En la revisión consultada de SonarQube, el control de calidad aparece como
superado. La captura de esta revisión se muestra en la figura~\ref{fig:aspectos-sonarqube}.

\begin{figure}[htbp]
    \centering
    \includegraphics[width=\textwidth]{memoria/aspectos-relevantes-del-desarrollo-del-proyecto/metricas-sonarqube.png}
    \caption{Métricas globales de calidad en SonarQube}
    \label{fig:aspectos-sonarqube}
\end{figure}

No se registran incidencias abiertas de seguridad ni de fiabilidad,
y ambas categorías obtienen una calificación A. La mantenibilidad también
obtiene una A, aunque conserva 43 incidencias abiertas. La cobertura alcanza
el 88,4\,\% sobre 8,4 mil líneas, mientras que la duplicación se sitúa en el
0,1\,\% sobre 35 mil líneas. No hay incidencias aceptadas y la sección de
\textit{Security Hotspots} no registra elementos.

Estas cifras describen el estado global del código en esa revisión y
complementan las comprobaciones ejecutadas en la integración continua. La
calificación no implica que no quede trabajo de mantenimiento, ya que las 43
incidencias abiertas corresponden precisamente a aspectos que pueden seguir
revisándose aunque el control de calidad se haya superado.

Para consultar qué cambia entre versiones, se mantiene manualmente el archivo
\texttt{CHANGELOG.md}, organizado por versión y tipo de cambio según
\textit{Keep a Changelog}~\cite{keep_a_changelog}. La numeración de las versiones
sigue la especificación SemVer, que distingue entre cambios incompatibles,
nuevas funcionalidades que mantienen la compatibilidad y correcciones de
errores que también la conservan~\cite{semantic_versioning}.

En GitHub, la publicación de una \textit{release} utiliza un flujo separado, activado
por un \textit{tag} de versión como \texttt{v1.0.0} o de forma manual. Antes
de publicarla, comprueba que su \textit{commit} pertenezca a la rama
principal, que las versiones declaradas coincidan y que exista su apartado
en \texttt{CHANGELOG.md}. Después prepara los archivos del código fuente
y sus sumas de verificación. Este flujo no vuelve a ejecutar las pruebas
ni depende del resultado de la \acs{CI}. En GitLab, la publicación se activa
con una etiqueta de versión y depende de varias comprobaciones previas,
entre ellas las pruebas del servidor y de la interfaz.

Para la documentación se utiliza despliegue continuo, \acs{CD}, mediante
Cloudflare Workers Builds, que conecta los cambios del repositorio con la
construcción y publicación del sitio~\cite{cloudflare_workers_builds}.
\acused{CD} Su flujo de GitHub Actions comprueba los archivos de Astro y
construye el sitio cuando cambian sus fuentes o su configuración, y también
admite ejecución manual. La aplicación se arranca a demanda mediante
Docker Compose, como se explica en el apartado~\ref{aspectos:publicacion}.

\section{Publicación de la aplicación}
\label{aspectos:publicacion}

La publicación se organiza en tres sitios con funciones distintas:

\begin{itemize}
    \item \href{https://manualito.dev/}{\textit{manualito.dev}} presenta el
    proyecto y sirve como punto de entrada.
    \item \href{https://app.manualito.dev/}{\textit{app.manualito.dev}}
    corresponde a la aplicación, cuyo despliegue se activa a demanda.
    \item \href{https://docs.manualito.dev/}{\textit{docs.manualito.dev}}
    expone las guías de uso y la documentación técnica, preparadas con
    Astro y Starlight~\cite{astro_documentation,starlight_documentation}.
\end{itemize}

Esta separación permite consultar los sitios estáticos, es decir, la presentación y la documentación sin
depender de que estén activos los servicios de procesamiento. La
aplicación, en cambio, necesita mantener disponibles los modelos, las
bases de datos y los procesos de trabajo, por lo que se ejecuta mediante
Docker Compose en un equipo con los recursos necesarios a modo de servidor.

El acceso exterior pasa por Cloudflare Tunnel, que establece una conexión
saliente desde ese equipo, y por Caddy, que sirve la interfaz y dirige las
peticiones al servidor~\cite{cloudflare_tunnel}. La configuración de
infraestructura se conserva en Terraform y los datos se guardan en
volúmenes persistentes para mantenerlos al recrear los contenedores. El
Anexo D recoge la preparación del entorno y las instrucciones de ejecución.

\section{Declaración del uso de inteligencia artificial generativa}
\label{aspectos:uso-ia}

Desarrollar un proyecto del tamaño de \textit{Manualito} ofrecía una
oportunidad para aprender sobre desarrollo agéntico~\cite{ibm_agentic_coding}, como se planteó en
los objetivos personales del apartado~\ref{objetivos:personales}. Además,
aplicar esta metodología en momentos puntuales, permite dedicar menos tiempo a tareas sencillas y repetitivas para poder
profundizar en las tecnologías utilizadas y concentrar el esfuerzo de
diseño y desarrollo en las partes que aportaban más valor a la aplicación.

Es por ello que, durante el desarrollo del proyecto se utilizaron Codex, de OpenAI,
y Claude Code, de Anthropic~\cite{claude_code_documentacion}, con los modelos
recogidos en la tabla~\ref{tabla:aspectos-modelos-ia}. Su ayuda se centró en
escribir código repetitivo y ayudar a buscar soluciones a problemas que dificultaban
avanzar, siempre a partir de instrucciones concretas y sin delegar las
decisiones sobre la arquitectura o el diseño del sistema.

Estas herramientas también ayudaron a elaborar y ejecutar pruebas una vez indicados
los respectivos casos de prueba, incluida
la interacción con el navegador mediante Playwright descrita en el
apartado~\ref{aspectos:pruebas} con la finalidad de obtener mejores resultados de las pruebas de
aceptación. Para valorar las soluciones propuestas se
recurrió a la documentación técnica y a la comprobación del comportamiento
de la aplicación, junto con las pruebas y las revisiones automáticas
recogidas en el apartado~\ref{aspectos:automatizacion}.

Durante la investigación también se recurrió a estas herramientas para
localizar artículos científicos y sintetizar la información relevante, con
el fin de facilitar la comparación de alternativas y la toma de decisiones.
Los resúmenes y las referencias propuestas se contrastaron con las
publicaciones originales. Además, estas herramientas ayudaron a escribir los datos sintéticos y
parte del código de los \textit{benchmarks} a partir de las indicaciones sobre qué
se quería medir y cómo hacerlo. En la evaluación de recuperación se generaron
20 de las 24 preguntas con un modelo, a las que se añadieron cuatro casos
procedentes de fallos detectados. En las comparaciones de lenguaje, otros
modelos actuaron como jueces mediante criterios escritos y se conservaron
sus desacuerdos~\cite{manualito_benchmarks}.

La interpretación de los resultados no se delegó enteramente en la clasificación obtenida, sino que, se valoraron
los resultados según las necesidades del proyecto, como reservar memoria
gráfica para la conversación o limitar la corrección del texto para no
introducir reglas falsas. Las simulaciones de uso real complementaron la
selección de los modelos, como se explica en el
apartado~\ref{aspectos:investigacion}.

En la memoria se ha recurrido a estas herramientas para mejorar la presentación
de los diagramas, ajustando su distribución, tipografía y apariencia para
facilitar la lectura sin cambiar el contenido técnico que se quería
representar. Todo ello a partir de instrucciones específicas.

Para precisar las versiones utilizadas, la tabla~\ref{tabla:aspectos-modelos-ia}
recoge las fechas de lanzamiento de los modelos considerados, incluidos
los que ya estaban disponibles al comenzar el proyecto. La fecha corresponde al
primer lanzamiento anunciado, aunque el acceso inicial fuera limitado.

\begin{table}[htbp]
    \centering
    \small
    \begin{tabular}{@{}p{0.30\textwidth}p{0.35\textwidth}p{0.27\textwidth}@{}}
        \toprule
        \textbf{Modelo} & \textbf{Fecha de lanzamiento} & \textbf{Laboratorio} \\
        \midrule
        Claude Sonnet 4.5 & 29/09/2025 & Anthropic~\cite{anthropic_sonnet45_lanzamiento} \\
        \addlinespace
        Claude Haiku 4.5 & 15/10/2025 & Anthropic~\cite{anthropic_haiku45_lanzamiento} \\
        \addlinespace
        Claude Sonnet 4.6 & 17/02/2026 & Anthropic~\cite{anthropic_sonnet46_lanzamiento} \\
        \addlinespace
        GPT-5.4 & 05/03/2026 & OpenAI~\cite{openai_gpt54_lanzamiento} \\
        \addlinespace
        GPT-5.5 & 23/04/2026 & OpenAI~\cite{openai_gpt55} \\
        \addlinespace
        GPT-5.6 Luna~\cite{openai_gpt56_luna} & 26/06/2026~\cite{openai_gpt56_previa} & OpenAI \\
        \addlinespace
        Claude Sonnet 5 & 30/06/2026 & Anthropic~\cite{anthropic_sonnet5_lanzamiento} \\
        \bottomrule
    \end{tabular}
    \caption{Modelos utilizados en algún momento del ciclo de desarrollo de \textit{Manualito}.
    Elaboración propia a partir de los anuncios y la documentación de los laboratorios.}
    \label{tabla:aspectos-modelos-ia}
\end{table}
\FloatBarrier

Para documentar esta ayuda se han tomado como referencia las orientaciones de
la Comisión Europea, que recomiendan identificar las herramientas
utilizadas, explicar su aportación y revisar sus resultados, manteniendo
la responsabilidad sobre el trabajo
presentado~\cite{comision_europea2026_ia_investigacion}.
